\section{Presentation-Perturbation Mini-Suite}
\label{app:presentation-perturbation-mini-suite}

The presentation-perturbation mini-suite froze 20 role-uncued source tasks from \code{eha-mvp/data/uncued-pilot-v1}: five each from \code{generated_lore}, \code{buried_primary}, \code{conflicting_evidence}, and \code{false_consensus}. Within each condition, the frozen slice used two \code{packet_judgment} tasks, two \code{evidence_selection} tasks, and one \code{active_verification} task where available. The suite used only the \code{neutral_metadata_visible} view and seed 20260525.

The run compared \code{baseline_original} against four paired perturbations: \code{order_randomized}, \code{source_type_masked}, \code{prompt_paraphrase}, and \code{citation_masked}. The schema and scorer were held fixed to the Phase 1 interface. Pre-model gates checked that the selection manifest referenced only the role-uncued pilot data, that prompt parity changed only the intended field for each perturbation, that hidden labels and condition or family labels were absent from model-visible prompts, that projected cost stayed below the hard cap, and that every perturbation had a paired baseline row.

The full run completed 200 planned calls: 20 source tasks times two models times five variants. All 200 latest rows parsed successfully. The cost report recorded USD 2.614767 spent against a USD 7 hard cap. The prompt audit and stored hidden-label audit both passed with zero hits. The paired analysis produced 160 baseline-versus-perturbation comparisons and 57 flip-review rows for local inspection.

The run directory contains the mini-suite result report, paired CSV tables, model-call artifacts, prompt audits, and the flip-review sheet. The top-level reports directory contains the Markdown result summary and the machine-readable JSON summary.
